\abstract{%
In this paper we evaluate multimodal models. Multimodal models are usually
evaluated by assigning them a task and then investigating whether
performance drops when one input is degraded. This mixes up two different
things: information that is lost due to degradation, and a model that still
keeps relying on an input after it has stopped helping. We propose a
black-box diagnostic that separates the two. Along a severity path we pair
utility with reliance and locate where utility falls below the no-input
level. We evaluate our measures on synthetic fusion models implemented with
controlled reliance, including an oracle control. On three open VLMs and a
multiple-choice VQA task, a swapped-image corruption leaves normalized
reliance at \checknum{0.487}{Qwen:swapblend:R_hat_s1} to
\checknum{0.578}{llava:swapblend:R_hat_s1} of its clean level,
\checknum{0.112}{paired:SmolVLM:delta_rhat} to
\checknum{0.319}{paired:Qwen:delta_rhat} above Gaussian noise,
and produces a significant utility inversion for one model under both
baselines. The pre-declared criterion was met in only
\checknum{2}{criterion:status:count_met} of the 42 combinations, and
conclusions for individual models depend on how the no-input baseline was
defined. A small modality-dropout study preserves clean accuracy but shows a
clear drop in raw reliance, while exposing a weakness of R(0)-normalized
measures.}
\keywords{Multimodal Learning, Vision-Language Models, Robustness Evaluation, Intelligent Systems}
